\section{Dynamique de Perron et mode stable de mémoire}

Sur un parcours fermé, le cocycle cylindrique de Perron satisfait

\begin{equation}
V_n=\varphi^n.
\label{rm:eq:cosmo-perron-V}
\end{equation}

En dimension spatiale \(D=3\), l’échelle déterminée par le cocycle est

\begin{equation}
\frac{a_n}{a_0}=V_n^{1/3}=\varphi^{n/3},
\label{rm:eq:cosmo-perron-a}
\end{equation}

tandis que le mode stable de l’action quadratique est

\begin{equation}
M_n=\varphi^{-2n}=V_n^{-2}.
\label{rm:eq:cosmo-perron-M}
\end{equation}

Éliminer \(n\) entre \cref{rm:eq:cosmo-perron-a,rm:eq:cosmo-perron-M} donne le résultat central :

\begin{theorem}[Loi cylindrique de sixième ordre pour la mémoire]
Pour tout niveau fermé de la réalisation de Perron dans \(D=3\),
\begin{equation}
\boxed{
M_n=\left(\frac{a_n}{a_0}\right)^{-6}.}
\label{rm:eq:cosmo-minus-six}
\end{equation}
L’exposant \(-6\) n’est pas ajusté : c’est le quotient de l’exposant stable \(-2\) par la racine dimensionnelle \(1/3\).
\end{theorem}

Aux retours distingués, on obtient

\[
\frac{a_9}{a_0}=\varphi^3,
\qquad
\frac{a_{27}}{a_0}=\varphi^9,
\qquad
\frac{a_{108}}{a_0}=\varphi^{36},
\]

et les mêmes niveaux vérifient exactement \(M_n=(a_n/a_0)^{-6}\). Si l’on déclare le paramètre temporel additif \(t_n=nt_0\), le taux logarithmique est constant :

\begin{equation}
H_\varphi
:=\frac{\Delta\log a_n}{\Delta t}
=\frac{\log\varphi}{3t_0},
\qquad
H_\varphi t_0
=0.1604039416865344824992529711\ldots.
\label{rm:eq:cosmo-Hphi}
\end{equation}

Le symbole \(H_\varphi\) désigne le taux de cette échelle discrète. Sa lecture comme paramètre de Hubble exige le transducteur cosmographique ; l’identité interne ne dépend pas de cette promotion.

\section{Réalisation dans le système d’Einstein--Cartan}

Dans un fluide de spin d’Einstein--Cartan, la densité quadratique de spin obéit à

\begin{equation}
s^2\propto a^{-6}.
\label{rm:eq:cosmo-EC-scaling}
\end{equation}

\cref{rm:eq:cosmo-minus-six} fournit l’exposant et le mode stable sans consulter un rebond. Pour effectuer la lecture physique, il faut déclarer une section positive \(s_*^2\) et une application

\begin{equation}
\mathcal T_{\rm EC}:M_n\longmapsto
s_n^2:=s_*^2M_n
=s_*^2\left(\frac{a_n}{a_0}\right)^{-6}.
\label{rm:eq:cosmo-EC-transducer}
\end{equation}

L’application conserve la loi d’échelle, mais ne détermine pas à elle seule l’amplitude \(s_*^2\), le courant de spin, le signe du terme de contorsion ni l’équation d’état. Avec une densité de masse, une équation effective prend la forme

\begin{equation}
H^2=\frac{8\pi G}{3}
\bigl(\rho_{\rm m}-\rho_{\rm spin,m}\bigr)+\cdots,
\label{rm:eq:cosmo-EC-mass}
\end{equation}

tandis qu’avec une densité d’énergie, il faut écrire

\begin{equation}
H^2=\frac{8\pi G}{3c^2}
\bigl(\rho_{\rm E}-\rho_{\rm spin,E}\bigr)+\cdots.
\label{rm:eq:cosmo-EC-energy}
\end{equation}

Les deux conventions ne se mélangent pas. Un candidat au rebond exige, outre l’égalité des densités, régularité, matière et données aux limites.

\begin{falsifier}
La transduction d’Einstein--Cartan échoue si la fonctionnelle d’échelle produit un exposant différent de \(-6\), si \(V_n\ne\varphi^n\) sur le parcours déclaré, ou si l’application physique ne conserve pas la positivité et les dimensions. Choisir \(s_*^2\) rétrospectivement pour imposer un seuil de rebond ne démontre pas le transducteur.
\end{falsifier}

\section{Comparaison structurelle entre le cycle de 108 états et la dynamique de Perron}

\begin{center}
\begin{tabularx}{\textwidth}{>{\bfseries}p{29mm}XX}
\toprule
& CT108--de Sitter & Perron--Einstein--Cartan\\
\midrule
Antécédent & retour circulaire d’ordre \(108\) & cocycle cylindrique et mode stable\\
Sélecteur & \(r_g=R_{108}\) & \(V_n=\varphi^n\), \(D=3\)\\
Magnitudes determinadas & rayon, courbure, température et entropie & expansion, taux et dilution de la mémoire\\
Invariant & \(G_a\hbar_a\) et rayon circulaire & \(M_na_n^6\)\\
Statut & interface de Sitter conditionnée & loi interne exacte ; lecture EC conditionnée\\
Critère de réfutation & rayon ou équation de fond incompatibles & perte de l’exposant \(-6\) ou du transducteur\\
\bottomrule
\end{tabularx}
\end{center}

\begin{statusbox}
Les identités CT108, le sélecteur \(\theta_{108}\) et la loi \(M_n=(a_n/a_0)^{-6}\) sont \status{Théorème interne exact}. La lecture du même rayon comme rayon de Sitter et l’identification de \(M_n\) à une densité de spin sont
\status{Réalisations physiques conditionnées}, chacune dotée de son propre critère de réfutation.
\end{statusbox}
